\chapter{Gérer un portefeuille multi-stratégies}\label{ch:s1:running-a-multi-strategy-book}

Pendant la première semaine d'août 2007, de nombreux fonds actions quantitatifs ont perdu d'un seul coup. Khandani et Lo, qui reconstituent l'épisode à partir de données de facteurs et de transactions,
mettent en évidence un désendettement généralisé du marché~: deux débouclages, le 1 août puis dès l'ouverture du 6 août, commençant par les valeurs financières, en position longue sur le ratio book-to-market et courte sur le
momentum des résultats, positions que de nombreux fonds détenaient sous des noms différents. Une société qui gère plusieurs stratégies est exposée au même risque au sein de
sa propre structure. Dans la société synthétique de dix \omterm{def:s1:running-a-multi-strategy-book:multi}{pods} de ce chapitre, cinq partagent des facteurs cachés qui s'effondrent ensemble trois fois en vingt ans~; détecter le chevauchement
et le plafonner réduit le drawdown maximal de la société de 43.5~\% à 35.7~\% pour une volatilité de 10~\%, et une \omterm{def:s1:running-a-multi-strategy-book:limit}{limite de drawdown} de 5~\% sur chaque \omterm{def:s1:running-a-multi-strategy-book:multi}{pod} arrête 91.6~\% des
\omterm{def:s1:running-a-multi-strategy-book:multi}{pods} sains dans l'année. La construction est \texttt{firm.multistrat}.

\section{Allocation entre stratégies}

\begin{definition}[Portefeuille multi-stratégies, pod]\label{def:s1:running-a-multi-strategy-book:multi}
Un \emph{portefeuille multi-stratégies}\index{portefeuille multi-stratégies} (multi-strategy book) est le portefeuille combiné d'une société regroupant plusieurs stratégies gérées par des équipes distinctes, auxquelles une fonction centrale alloue du capital ou
du risque, et elle compense leurs transactions et contrôle leur risque global. Un \emph{pod}\index{pod} est une équipe de gestion autonome, dotée de son propre capital, de ses
propres limites de risque et de son P\&L.
\end{definition}

\begin{definition}[Corrélation entre stratégies]\label{def:s1:running-a-multi-strategy-book:correlation}
La \emph{corrélation entre stratégies}\index{corrélation entre stratégies} (strategy correlation) est la corrélation entre les flux de P\&L de deux stratégies~; mesurée sur des périodes normales, elle peut
être faible alors que les stratégies partagent une exposition qui domine en période de tensions, si bien qu'une société l'estime à la fois à partir du P\&L et à partir des expositions aux facteurs des positions.
\end{definition}

La première décision de la société porte sur le risque alloué à chaque \omterm{def:s1:running-a-multi-strategy-book:multi}{pod}. La budgétisation du risque (livre~7, chapitre~26) donne la règle par défaut~: des poids inversement proportionnels à la volatilité glissante de chaque
\omterm{def:s1:running-a-multi-strategy-book:multi}{pod}, de sorte que chacun contribue à un risque comparable, ces poids étant rééquilibrés chaque mois (le \cref{lst:s1:running-a-multi-strategy-book:firm}). La société synthétique
compte dix \omterm{def:s1:running-a-multi-strategy-book:multi}{pods} à 10~\% de volatilité chacun, avec des ratios de Sharpe compris entre 0.5 et 1.0. Les \omterm{def:s1:running-a-multi-strategy-book:multi}{pods} 1 à 3 ont une charge de 0.4 sur un facteur commun, les \omterm{def:s1:running-a-multi-strategy-book:multi}{pods} 4 et 5 sur un autre~; les
facteurs sont aussi discrets que n'importe quel autre bruit, sauf lors de trois krachs de dix jours, chacun représentant une chute de six volatilités de \omterm{def:s1:running-a-multi-strategy-book:multi}{pod}. L'edge du \omterm{def:s1:running-a-multi-strategy-book:multi}{pod} 10 s'éteint après dix ans.

\section{Compensation et risque partagé}

Mesurée sur vingt ans de P\&L quotidiens, la corrélation moyenne entre les trois \omterm{def:s1:running-a-multi-strategy-book:multi}{pods} du premier facteur est de 0.22~; celle d'un \omterm{def:s1:running-a-multi-strategy-book:multi}{pod} exposé à aucun des deux facteurs, $-0.004$
avec le premier. La corrélation sous-estime ce qu'ils partagent. Lors de chaque krach, les trois \omterm{def:s1:running-a-multi-strategy-book:multi}{pods} perdent ensemble 2.4 à 2.6 de leurs volatilités annuelles, soit 24~\% à
26~\% du capital, en dix jours. La même position prise sous des noms différents est invisible dans les corrélations mesurées sur des années calmes et évidente dans les
positions des \omterm{def:s1:running-a-multi-strategy-book:multi}{pods}~: un modèle de risque (livre~7, chapitre~24) projette les avoirs de chaque \omterm{def:s1:running-a-multi-strategy-book:multi}{pod} sur les facteurs, et la société peut appliquer un cap à l'exposition totale à chacun.

\begin{center}
\small
\begin{tabular}{@{}lccc@{}}
\toprule
société à 10~\% de volatilité & sans contrôle du chevauchement & cap de facteur 1.0 & cap de facteur 0.5\\
\midrule
ratio de Sharpe & 1.45 & 1.52 & 1.70\\
drawdown maximal & 43.5~\% & 42.8~\% & 35.7~\%\\
perte lors des trois krachs & 31.5, 31.7, 33.6~\% & 30.1, 30.1, 32.3~\% & 21.6, 20.7, 22.9~\%\\
\bottomrule
\end{tabular}
\end{center}

Le cap limite l'exposition totale de chaque facteur commun à l'équivalent d'un \omterm{def:s1:running-a-multi-strategy-book:multi}{pod} (ou d'un \omterm{def:s1:running-a-multi-strategy-book:multi}{demi-pod}) à pondération égale, en réduisant les \omterm{def:s1:running-a-multi-strategy-book:multi}{pods} qui le portent. Plafonner à
un \omterm{def:s1:running-a-multi-strategy-book:multi}{demi-pod} réduit chaque krach d'environ un tiers et relève le ratio de Sharpe de la société de 1.45 à 1.70 (la \cref{fig:s1:running-a-multi-strategy-book:paths}),
car les krachs représentaient le principal coût des \omterm{def:s1:running-a-multi-strategy-book:multi}{pods} qui se chevauchent. Les pertes sont importantes parce qu'une société de dix \omterm{def:s1:running-a-multi-strategy-book:multi}{pods} presque décorrélés doit leur appliquer un effet de levier pour
atteindre une volatilité de 10~\%, et l'effet de levier amplifie tout ce qu'ils partagent.

\begin{omfigure}
\begin{tikzpicture}
\begin{axis}[width=11cm,height=5.2cm,xlabel={\small année},ylabel={\small rendement cumulé (\%)},tick label style={font=\footnotesize},
  xmin=1,xmax=20,ymin=-20,ymax=340,grid=major,grid style={gray!25},legend pos=north west,legend style={font=\scriptsize},table/col sep=comma]
\fill[omExo!15] (axis cs:5.952,-20) rectangle (axis cs:5.992,340);
\fill[omExo!15] (axis cs:14.286,-20) rectangle (axis cs:14.325,340);
\fill[omExo!15] (axis cs:17.460,-20) rectangle (axis cs:17.500,340);
\addplot[omDef,thick] table[x=year,y=none]{figdata/strategies-1/28-running-a-multi-strategy-book/paths.csv};
\addplot[omThm,thick] table[x=year,y=cap]{figdata/strategies-1/28-running-a-multi-strategy-book/paths.csv};
\legend{sans contrôle du chevauchement,facteurs communs plafonnés à un demi-pod}
\end{axis}
\end{tikzpicture}

\omcaption{La société synthétique de dix \omterm{def:s1:running-a-multi-strategy-book:multi}{pods}, calibrée à 10~\% de volatilité, avec et sans plafond sur les expositions des \omterm{def:s1:running-a-multi-strategy-book:multi}{pods} aux facteurs communs~; sommes cumulées des
rendements quotidiens sur vingt ans, les bandes ombrées marquant les trois krachs des facteurs communs. Données~: \texttt{s1\_multistrat.paths}.}\label{fig:s1:running-a-multi-strategy-book:paths}
\end{omfigure}

La même vue centralisée permet aussi de réduire les coûts. Lorsque des \omterm{def:s1:running-a-multi-strategy-book:multi}{pods} négocient les mêmes instruments, la société n'exécute que le solde net~: le croisement interne (livre~7, chapitre~27). Dix
\omterm{def:s1:running-a-multi-strategy-book:multi}{pods} négociant 200 instruments de façon indépendante, chacun en négociant un cinquième, voient 34.6~\% de leur volume brut se compenser en interne~; si tous négocient tous
les instruments, 68.5~\%, proche de $1 - 1/\sqrt{10}$.

\section{Désactiver les stratégies}

\begin{definition}[Limite de drawdown]\label{def:s1:running-a-multi-strategy-book:limit}
Une \emph{limite de drawdown}\index{limite de drawdown} (drawdown limit) arrête ou réduit le risque d'une stratégie lorsque sa perte depuis son sommet précédent dépasse un seuil, une fraction de son
capital alloué~; elle arbitre entre le coût d'arrêter de bonnes stratégies victimes de la malchance et celui de faire tourner trop longtemps des stratégies mortes.
\end{definition}

Toute règle d'arrêt commet deux erreurs. Elle arrête des \omterm{def:s1:running-a-multi-strategy-book:multi}{pods} dont l'edge est encore vivant, et laisse tourner les \omterm{def:s1:running-a-multi-strategy-book:multi}{pods} morts. La \Cref{fig:s1:running-a-multi-strategy-book:stops} montre
la première~: un \omterm{def:s1:running-a-multi-strategy-book:multi}{pod} de ratio de Sharpe 0.7 à 10~\% de volatilité atteint un drawdown de 5~\% dans l'année sur 91.6~\% des 20\,000 trajectoires simulées, un drawdown de 10~\% sur
34.5~\%, 20~\% sur 2.1~\%. Une limite se mesure en volatilités propres au \omterm{def:s1:running-a-multi-strategy-book:multi}{pod}~: 5~\% équivaut à la moitié de l'écart-type annuel, bien dans le bruit ordinaire. Dans la
société synthétique, où les \omterm{def:s1:running-a-multi-strategy-book:multi}{pods} arrêtés redémarrent chaque année, une limite de 5~\% arrête 85.8 \omterm{def:s1:running-a-multi-strategy-book:multi}{pods} vivants pour cent \omterm{def:s1:running-a-multi-strategy-book:multi}{années-pod}~; à 20~\%, elle en arrête encore 14.7,
car les krachs communs font franchir ensemble cette limite à plusieurs \omterm{def:s1:running-a-multi-strategy-book:multi}{pods}.

\begin{omfigure}
\begin{tikzpicture}
\begin{axis}[width=10cm,height=5cm,xlabel={\small limite de drawdown (\% du capital du pod)},ylabel={\small \% des pods vivants arrêtés dans l'année},
  tick label style={font=\footnotesize},xmin=3,xmax=32,ymin=0,ymax=100,grid=major,grid style={gray!25},
  legend style={font=\scriptsize,at={(0.5,-0.3)},anchor=north},table/col sep=comma]
\addplot[omThm,thick,mark=*] table[x=limit_pct,y=paths_pct]{figdata/strategies-1/28-running-a-multi-strategy-book/stops.csv};
\addplot[omDef,thick,mark=square*] table[x=limit_pct,y=firm_pct]{figdata/strategies-1/28-running-a-multi-strategy-book/stops.csv};
\legend{un pod (ratio de Sharpe 0.7) sur 20\,000 trajectoires,pods de la société synthétique}
\end{axis}
\end{tikzpicture}

\omcaption{Faux arrêts~: la part des \omterm{def:s1:running-a-multi-strategy-book:multi}{pods} dont l'edge est vivant et qui ont atteint une \omterm{def:s1:running-a-multi-strategy-book:limit}{limite de drawdown} dans l'année, pour un \omterm{def:s1:running-a-multi-strategy-book:multi}{pod} isolé de ratio de Sharpe 0.7 à 10~\%
volatilité (trajectoires simulées) et pour les \omterm{def:s1:running-a-multi-strategy-book:multi}{pods} de la société synthétique, dont les krachs communs ajoutent des arrêts. Données~: \texttt{s1\_multistrat.limits}.}\label{fig:s1:running-a-multi-strategy-book:stops}
\end{omfigure}

La seconde erreur est le \omterm{def:s1:running-a-multi-strategy-book:multi}{pod} mort. Son ratio de Sharpe passe de 0.9 à $-0.3$ après dix ans~; la limite de 5~\% l'arrête 98 jours de bourse plus tard, la limite de 10~\%
après 307, la limite de 15~\% après 979, et les limites plus larges ne l'arrêtent jamais, sur les dix ans restants. Une limite serrée attrape vite le \omterm{def:s1:running-a-multi-strategy-book:multi}{pod} mort en attrapant presque chaque
\omterm{def:s1:running-a-multi-strategy-book:multi}{pod} rapidement. Le ratio de Sharpe de la société selon les limites (1.49 à 5~\%, 1.74 à 10~\%, 1.61 à 15~\%, 1.48 à 20~\%) culmine au milieu, et sur un échantillon de cette taille
ce pic n'est pas un optimum robuste.

\section{Capital, limites de drawdown et modèle de pod}

La \omterm{def:s1:running-a-multi-strategy-book:limit}{limite de drawdown} repose sur une théorie. Grossman et Zhou ont résolu le problème d'un investisseur qui ne veut jamais voir sa richesse passer sous une fraction fixe de son
maximum courant~: la position risquée optimale est proportionnelle à la distance à ce plancher. Le risque baisse à mesure que les pertes approchent du plancher, au lieu de basculer
à zéro au plancher, et il remonte à mesure que le surplus se reconstitue. Un modèle de \omterm{def:s1:running-a-multi-strategy-book:multi}{pod} qui réduit de moitié le risque d'un \omterm{def:s1:running-a-multi-strategy-book:multi}{pod} à un premier seuil et l'arrête à un second est une version à deux étapes
de cette règle. Les enseignements centraux du chapitre en découlent. Les limites doivent s'exprimer en unités de volatilité de chaque \omterm{def:s1:running-a-multi-strategy-book:multi}{pod} et de ratio de Sharpe attendu, et non en
pourcentages ronds. Les expositions communes doivent être mesurées à partir des positions, et non des corrélations de P\&L. Et le capital de la société est le plus sûr lorsque les drawdowns sont
contrôlés en continu, comme le fait la règle de Grossman et Zhou, plutôt que par une rupture brutale qui arrête de nombreux bons \omterm{def:s1:running-a-multi-strategy-book:multi}{pods} pour en attraper quelques mauvais.

\section{Fiches de stratégie}

\begin{strategyfile}{Allocation multi-stratégies avec contrôle des recouvrements}\label{strat:s1:running-a-multi-strategy-book:overlap}
\sfield{Qui vous paie, et pourquoi} Les edges combinés des \omterm{def:s1:running-a-multi-strategy-book:multi}{pods}, avec la diversification lorsque leurs risques sont réellement différents.
\sfield{Instruments et plateformes} Tout ce que tradent les \omterm{def:s1:running-a-multi-strategy-book:multi}{pods}.
\sfield{Signal} La volatilité glissante des \omterm{def:s1:running-a-multi-strategy-book:multi}{pods} et les expositions aux facteurs de leurs positions.
\sfield{Dimensionnement et exécution} Budgets de risque en inverse de la volatilité~; caps sur chaque facteur partagé~; transactions compensées entre \omterm{def:s1:running-a-multi-strategy-book:multi}{pods}.
\sfield{Coûts} Plus faibles grâce à la compensation~; un modèle de risque et un desk d'exécution centralisé.
\sfield{Comment il meurt} Des facteurs encombrés, partagés par de nombreuses sociétés, comme en août 2007.
\sfield{Horizon, capacité, infrastructure} Continu~; données de position de chaque \omterm{def:s1:running-a-multi-strategy-book:multi}{pod}.
\sfield{Backtest honnête} Expositions issues des positions à chaque date~; épisodes de krach dans l'échantillon.
\sfield{Sources} Khandani et Lo (2011)~; ce chapitre~: drawdown de 43.5~\% à 35.7~\% avec un cap d'une \omterm{def:s1:running-a-multi-strategy-book:multi}{demi-pod}.
\end{strategyfile}

\begin{strategyfile}{Exposition contrôlée par le drawdown}\label{strat:s1:running-a-multi-strategy-book:drawdown}
\sfield{Qui vous paie, et pourquoi} Personne~: c'est une règle de risque qui échange du rendement espéré contre un plancher sur les pertes.
\sfield{Instruments et plateformes} Toute stratégie ou la société.
\sfield{Signal} La distance entre le plus haut courant et un plancher.
\sfield{Dimensionnement et exécution} Risque proportionnel au surplus au-dessus du plancher, ou réduit par paliers.
\sfield{Coûts} Les faux stops et le rendement abandonné après eux.
\sfield{Comment il meurt} Des limites fixées en pourcentages plutôt qu'en volatilités~; des krachs communs qui frappent de nombreux \omterm{def:s1:running-a-multi-strategy-book:multi}{pods} ensemble.
\sfield{Horizon, capacité, infrastructure} Continu~; P\&L quotidien.
\sfield{Backtest honnête} Comptabiliser les faux stops autant que les pertes évitées.
\sfield{Sources} Grossman et Zhou (1993)~; ce chapitre~: 91.6~\% des \omterm{def:s1:running-a-multi-strategy-book:multi}{pods} sains arrêtés en un an à une limite de 5~\%.
\end{strategyfile}

\section{Tutoriel~: cinq à la fois}\label{tut:s1:running-a-multi-strategy-book:tut}

\begin{tutorial}
\textbf{Objectif.} Faire tourner une société de dix \omterm{def:s1:running-a-multi-strategy-book:multi}{pods} aux facteurs partagés cachés, allouer avec et sans contrôle des recouvrements, appliquer des \omterm{def:s1:running-a-multi-strategy-book:limit}{limites de drawdown} et mesurer les faux
stops et la survie du \omterm{def:s1:running-a-multi-strategy-book:multi}{pod} mort, et compenser les transactions. \textbf{État final~:} les tableaux et les deux figures.
\begin{tutsteps}
\item \textbf{Allocation et société}~: budgets de risque, le cap sur les facteurs, stops de drawdown.
\omcode{code/firm/multistrat/firm_multistrat.py}{59}{95}{Budgets de risque avec cap sur les facteurs, et société avec limites de drawdown.}\label{lst:s1:running-a-multi-strategy-book:firm}
\item \textbf{Faux stops et compensation}.
\omcode{code/firm/multistrat/firm_multistrat.py}{98}{110}{La probabilité d'un faux stop, et la compensation.}\label{lst:s1:running-a-multi-strategy-book:stops}
\item \textbf{Exécuter} \texttt{correlations()}, \texttt{overlap(cap)}, \texttt{limits(limit)}, \texttt{net\_savings()} et \texttt{fig\_multistrat.py}.
\end{tutsteps}
\textbf{Pour aller plus loin.} Remplacer la rupture par la règle proportionnelle de Grossman et Zhou~; estimer les facteurs partagés à partir du seul P\&L et voir
à quel point ils sont découverts tard~; intégrer le P\&L des stratégies propres à ce livre (chapitres 2, 5, 19 et 20) comme \omterm{def:s1:running-a-multi-strategy-book:multi}{pods}.
\end{tutorial}

\section{Construction~: la société multi-stratégies}\label{bld:s1:running-a-multi-strategy-book:multistrat}

\begin{build}
\bfield{Rôle} \omterm{def:s1:running-a-multi-strategy-book:multi}{Pods} simulés aux facteurs partagés, allocation par budgets de risque avec contrôle des recouvrements, \omterm{def:s1:running-a-multi-strategy-book:limit}{limites de drawdown}, taux de faux stops et compensation.
\bfield{Interface} \texttt{PodConfig(\dots)}, \texttt{simulate\_pods(cfg, rng)}, \texttt{allocate(pnl, load, t, window, factor\_cap)},
\texttt{run\_firm(S, rebalance, factor\_cap, limit)}, \texttt{stop\_rate(sr, vol, limit, years, n, rng)}, \texttt{netting(trades)}.
\bfield{Règles} Allocations fondées sur les données passées~; expositions calculées à partir des positions~; \omterm{def:s1:running-a-multi-strategy-book:multi}{pods} arrêtés redémarrés chaque année.
\bfield{Tests d'acceptation} \texttt{code/firm/multistrat/tests/}~: les facteurs partagés corrèlent les \omterm{def:s1:running-a-multi-strategy-book:multi}{pods} et les font chuter ensemble, le cap réduit le poids des \omterm{def:s1:running-a-multi-strategy-book:multi}{pods}
fortement exposés, la compensation vérifiée à la main, les faux stops qui diminuent quand la limite augmente.
\bfield{Pour aller plus loin} Contrôle proportionnel du drawdown~; réallocation du capital après les stops~; les stratégies propres au livre comme \omterm{def:s1:running-a-multi-strategy-book:multi}{pods}.
\end{build}

\begin{omsources}
\item A.~E.~Khandani et A.~W.~Lo, ``What happened to the quants in August 2007? Evidence from factors and transactions data'', \emph{Journal of
Financial Markets} 14(1), 2011.
\item S.~J.~Grossman et Z.~Zhou, ``Optimal investment strategies for controlling drawdowns'', \emph{Mathematical Finance} 3(3), 1993.
\item L.~H.~Pedersen, \emph{Efficiently Inefficient: How Smart Money Invests and Market Prices Are Determined}, Princeton University Press, 2015.
\end{omsources}

\section{Exercices}

\begin{exercise}[$\star$]\label{exo:s1:running-a-multi-strategy-book:1}
Dix \omterm{def:s1:running-a-multi-strategy-book:multi}{pods} tradent les mêmes instruments avec des transactions indépendantes de taille égale. Quelle part de leur volume brut la compensation permet-elle d'économiser~?
\end{exercise}

\begin{exercise}[$\star$]\label{exo:s1:running-a-multi-strategy-book:2}
Un \omterm{def:s1:running-a-multi-strategy-book:multi}{pod} charge 0.4 sur un facteur qui chute de six volatilités de \omterm{def:s1:running-a-multi-strategy-book:multi}{pod} lors d'un krach. Que perd le \omterm{def:s1:running-a-multi-strategy-book:multi}{pod}, en volatilités et en pourcentage, pour une volatilité de 10~\%~?
\end{exercise}

\begin{exercise}[$\star$]\label{exo:s1:running-a-multi-strategy-book:3}
Pourquoi trois \omterm{def:s1:running-a-multi-strategy-book:multi}{pods} de corrélation 0.22 perdent-ils ensemble lors d'un krach~?
\end{exercise}

\begin{exercise}[$\star\star$]\label{exo:s1:running-a-multi-strategy-book:4}
Une \omterm{def:s1:running-a-multi-strategy-book:limit}{limite de drawdown} de 5~\% arrête 91.6~\% des \omterm{def:s1:running-a-multi-strategy-book:multi}{pods} à une volatilité de 10~\% en un an. Quelle part arrête-t-elle pour des \omterm{def:s1:running-a-multi-strategy-book:multi}{pods} gérés à une volatilité de 5~\%, et pourquoi~?
\end{exercise}

\begin{exercise}[$\star\star$]\label{exo:s1:running-a-multi-strategy-book:5}
Pourquoi une limite serrée détecte-t-elle le \omterm{def:s1:running-a-multi-strategy-book:multi}{pod} mort le plus vite, et à quel prix~?
\end{exercise}

\begin{exercise}[$\star\star$]\label{exo:s1:running-a-multi-strategy-book:6}
Expliquer la règle de Grossman et Zhou et comment une limite de \omterm{def:s1:running-a-multi-strategy-book:multi}{pod} à deux étapes en fournit une approximation.
\end{exercise}

\begin{exercise}[$\star\star\star$]\label{exo:s1:running-a-multi-strategy-book:7}
\emph{Programmation.} Exécuter \texttt{overlap(0.25)} avec un cap plus strict. Que deviennent les krachs et le ratio de Sharpe, et qu'est-ce qui devrait inquiéter dans une
société réelle~?
\end{exercise}

\begin{exercise}[$\star\star\star$]\label{exo:s1:running-a-multi-strategy-book:8}
\emph{Trouver la faille.} «~Nos dix \omterm{def:s1:running-a-multi-strategy-book:multi}{pods} ont une corrélation moyenne deux à deux de 0.05 sur trois ans, donc le risque de la société est diversifié.~»
\end{exercise}

\section{Problème~: cinq à la fois}

\begin{problem}[{Problème du week-end --- une société de pods}]\label{pb:s1:running-a-multi-strategy-book:1}
La société synthétique du chapitre et le bilan public.

\medskip\noindent\textbf{Partie I --- La société.}
\begin{enumerate}
\item Définir un \omterm{def:s1:running-a-multi-strategy-book:multi}{portefeuille multi-stratégies}, un \omterm{def:s1:running-a-multi-strategy-book:multi}{pod} et la \omterm{def:s1:running-a-multi-strategy-book:correlation}{corrélation entre stratégies}.
\item Décrire la société synthétique~: \omterm{def:s1:running-a-multi-strategy-book:multi}{pods}, facteurs partagés, krachs et \omterm{def:s1:running-a-multi-strategy-book:multi}{pod} mort.
\item Comment la société alloue-t-elle le risque~?
\item Que Khandani et Lo ont-ils trouvé à propos d'août 2007~?
\end{enumerate}

\medskip\noindent\textbf{Partie II --- Risque partagé.}
\begin{enumerate}[resume]
\item Donner les corrélations mesurées sur vingt ans et les pertes des \omterm{def:s1:running-a-multi-strategy-book:multi}{pods} lors des krachs.
\item Comment la société détecte-t-elle le recouvrement~?
\item Donner le ratio de Sharpe de la société, le plus grand drawdown et les pertes lors des krachs avec et sans cap.
\item Quelle part la compensation permet-elle d'économiser~?
\end{enumerate}

\medskip\noindent\textbf{Partie III --- Stops.}
\begin{enumerate}[resume]
\item Définir une \omterm{def:s1:running-a-multi-strategy-book:limit}{limite de drawdown} et ses deux types d'erreurs.
\item Donner les taux de faux stops sur les trajectoires simulées et dans la société.
\item Combien de temps le \omterm{def:s1:running-a-multi-strategy-book:multi}{pod} mort reste-t-il en activité sous chaque limite~?
\item Pourquoi la meilleure limite de la société n'est-elle pas un optimum robuste~?
\end{enumerate}

\medskip\noindent\textbf{Partie IV --- Le verdict.}
\begin{enumerate}[resume]
\item Énoncer le \emph{résultat nommé}~: le drawdown de la société avec et sans détection du recouvrement, et le taux de faux stops d'une \omterm{def:s1:running-a-multi-strategy-book:limit}{limite de drawdown} de 5~\%.
\item Que Grossman et Zhou ont-ils montré~?
\item En quelles unités les limites doivent-elles être fixées~?
\item Pourquoi les krachs communs font-ils augmenter les faux stops pour les limites larges~?
\item Quelles données chaque \omterm{def:s1:running-a-multi-strategy-book:multi}{pod} doit-il transmettre au centre~?
\item Quel lien ce chapitre entretient-il avec le livre~7, chapitre~28~?
\item Quels \omterm{def:s1:running-a-multi-strategy-book:multi}{pods} couperiez-vous en premier après un krach, et pourquoi pas tous~?
\item En une phrase~: qu'apporte le centre d'une société multi-stratégies~?
\end{enumerate}
\end{problem}

\section{Questions d'entretien}

\begin{interviewq}[$\star$ \iqroles{risque}]\label{iq:s1:running-a-multi-strategy-book:1}
Comment fixeriez-vous une \omterm{def:s1:running-a-multi-strategy-book:limit}{limite de drawdown} pour un nouveau \omterm{def:s1:running-a-multi-strategy-book:multi}{pod}~?
\end{interviewq}

\begin{interviewq}[$\star\star$ \iqroles{risque}]\label{iq:s1:running-a-multi-strategy-book:2}
Deux \omterm{def:s1:running-a-multi-strategy-book:multi}{pods} sur des classes d'actifs différentes ont perdu 8~\% chacun la même semaine. Qu'est-ce qui retiendrait votre attention~?
\end{interviewq}

\begin{interviewq}[$\star\star$ \iqroles{chercheur}]\label{iq:s1:running-a-multi-strategy-book:3}
Comment estimeriez-vous si l'edge d'une stratégie a disparu, et combien de temps cela prendrait-il~?
\end{interviewq}

\begin{interviewq}[$\star\star$ \iqroles{développeur}]\label{iq:s1:running-a-multi-strategy-book:4}
Concevez le netting des transactions entre \omterm{def:s1:running-a-multi-strategy-book:multi}{pods} qui négocient les mêmes instruments.
\end{interviewq}

\begin{interviewq}[$\star\star$ \iqroles{trader}]\label{iq:s1:running-a-multi-strategy-book:5}
Votre \omterm{def:s1:running-a-multi-strategy-book:multi}{pod} est 4~\% sous son plus haut et la limite est de 5~\%. Que faites-vous, et que la société souhaite-t-elle que vous fassiez~?
\end{interviewq}

\begin{interviewq}[$\star\star\star$ \iqroles{chercheur}]\label{iq:s1:running-a-multi-strategy-book:6}
Le P\&L cumulé d'un \omterm{def:s1:running-a-multi-strategy-book:multi}{pod} est un mouvement brownien avec dérive $\mu$ et volatilité $\sigma$. En utilisant le fait que, pour $\mu = 0$, le drawdown par rapport au maximum courant
à l'instant $T$ suit la loi de $|W_T|$, donnez la probabilité que le drawdown dépasse $L$ à un instant $T$ donné, et
expliquez pourquoi la probabilité qu'il dépasse $L$ à un moment quelconque jusqu'à $T$ est plus grande.
\end{interviewq}
